\section{Mathematical Derivations and Model Tests}
\label{app:math}

\subsection{Conditional continuation tails and the novel-continuation rate}
For a context $c_i$ with $F_i=F$, let $y_{i,(r)}$ denote the $r$-th most likely continuation and $p_{i,(r)}$ its true conditional probability. Model the observed continuations $Y_{i,1},\ldots,Y_{i,F}$ as independent draws, with counts $n_{i,(r)}(F)=\#\{k\in\{1,\ldots,F\}:Y_{i,k}=y_{i,(r)}\}$. The novel-continuation rate $f_i(F)$ and its expectation under this sampling model are
\begin{equation}
f_i(F)=\sum_r p_{i,(r)}
\mathbf{1}\left[n_{i,(r)}(F)=0\right],
\quad
\mathbb{E}[f_i(F)]=\sum_r p_{i,(r)}(1-p_{i,(r)})^F.
\end{equation}

Assume a conditional tail $p_{i,(r)}\sim A_i r^{-s}$ with $s>1$. For large $F$ within the tail's scaling range, before finite-vocabulary effects become significant, the dominant probabilities are of order $1/F$. Approximating $(1-p_{i,(r)})^F$ by $e^{-Fp_{i,(r)}}$ and the rank sum by an integral then gives
\citep{gnedin2007occupancy}
\begin{equation}
\begin{aligned}
\mathbb{E}[f_i(F)]
&\sim
\int_1^\infty A_i r^{-s}e^{-F A_i r^{-s}}\,\mathrm{d}r\\
&\sim
\frac{A_i^{1/s}}{s}
\Gamma\left(1-\frac{1}{s}\right)
F^{-(1-1/s)}\\
&\propto F^{-\alpha},
\quad
\alpha=1-\frac{1}{s}.
\end{aligned}
\label{eq:occupancy}
\end{equation}

This provides a conditional explanation of the empirical scaling in Section~\ref{sec:coverage}. Finite vocabulary limits the power-law regime to an intermediate range. Heterogeneity across contexts may alter the observed effective exponent.

\subsection{Crowding-out within a context subspace}
\label{app:crowding-derivation}

Fix a context $c_i$ and call the next-token logits computed from it, $a_v$ for $v$ in the vocabulary, its \emph{context subspace}; the $n$-gram table gives $c_i$ a learned coordinate, so these logits can move without moving those of other contexts (Section~\ref{sec:context-locality}). Partition the vocabulary into the continuations of $c_i$ seen in training ($\mathcal{S}$), those that occur only in validation ($\mathcal{V}$), and the rest. With $p_v=\softmax{a_v}$, one idealized gradient step on a training target $C\in\mathcal{S}$ moves every logit by
\begin{equation}
a_v^{(t+1)}-a_v^{(t)}=-\eta\frac{\partial\ell_C}{\partial a_v}=-\eta\left(p_v^{(t)}-\mathbf{1}[v=C]\right),\qquad \eta>0,
\label{eq:av}
\end{equation}
raising $a_C$ by $\eta(1-p_C)$ and lowering every other logit by $\eta p_v$. For any $u\notin\mathcal{S}$ the odds $z(t)=p_C^{(t)}/p_u^{(t)}$ therefore satisfy
\begin{equation}
z(t+1)=z(t)\exp\!\left[\eta\left(1-p_C^{(t)}+p_u^{(t)}\right)\right]>z(t),
\label{eq:zt}
\end{equation}
and nothing saturates while $p_C<1$. When $\mathcal{S}$ is a single token, the training loss is $-\ln p_C$ and each pass transfers mass from $\mathcal{V}$ and the rest onto $C$; the validation loss on $\mathcal{V}$ rises by the same logit shift that lowers the training loss. Section~\ref{sec:sharpening} identifies $\ln z$, averaged over $\mathcal{V}$, with the depth $D_j$ of Equation~\eqref{eq:gap-decomp}.

\subsection{Position-class decomposition of the excess validation loss}
\label{app:novel-share}

Table~\ref{tab:novel-share} decomposes the both-table model's excess validation loss over the no-gram control (runs \code{s1v5_128_marm_both_ckpt_fixed} and \code{s1v5_128_marm_nogram_ckpt_fixed}, seed 42, step 1000) by whether the true continuation followed its context in the training shard. Contexts are counted at bigram or trigram granularity; \emph{novel} includes positions whose context itself is absent from training.

\begin{table}[!htbp]
\centering\small
\caption{Share of the step-1000 excess validation loss ($1.68$ nats per token, 589{,}248 positions) carried by each position class.}
\label{tab:novel-share}
\begin{tabular}{llrr}
\toprule
context & class & share of positions & share of excess loss \\
\midrule
bigram & novel continuation & 31\% & 73\% \\
bigram & novel or seen with $F\le 31$ & 39\% & 77\% \\
trigram & novel continuation & 66\% & 100\% \\
\bottomrule
\end{tabular}
\end{table}

%\subsection{Data volume, epoch length, and table capacity}
%\label{app:data-volume}

%If the context set is fixed, each $F(i)$ is multiplied by $k$, validation weights remain unchanged, and the per-event damage amplitude is fixed, a missing-mass law with a common exponent gives
%\begin{equation}
% \sum_i w_i A_j c_i[kF(i)]^{-\alpha} =k^{-\alpha}\sum_i w_i A_j c_iF(i)^{-\alpha}.
%\end{equation}
%These assumptions are restrictive. Adding independent data can change the context set, continuation coverage, and the effective frequency distribution. At a fixed step budget, it also changes the number of replays. Fixed-step and fixed-pass comparisons are therefore needed to separate coverage from repeated reinforcement.

%The measured decay is faster than coverage alone: in the dose sweep of Figure~\ref{fig:corpus-size}, the epoch-3 gap falls with corpus size with a local log-log slope of about $0.15$--$0.26$ below $2D_1$ but $0.8$--$0.95$ above it, whereas $k^{-\alpha}$ with $\alpha\in[0.23,0.38]$ predicts a constant slope of at most $0.38$. The missing factor is load suppression. The number of distinct contexts $K(D)$ grows with the corpus while $R$ stays fixed, and an empirical dilution surface built from the single-table bigram table-size sweep (\code{s1v5_128_tbl_bi1_R*_fixed}, seed 42, step 1000) shows a suppression factor $a(K/R)$ that is nearly flat in context frequency and scales approximately as $(K/R)^{-0.59}$ over loads $10$--$1000$; the same exponent accounts for the table-size slope of Figure~\ref{fig:capacity}. Combining the two factors, $\operatorname{gap}(D)\approx A(\text{passes})\cdot a(K(D)/R)\cdot\sum_f \mu_f M(f)$: coverage supplies the $D^{-\alpha}$ part, load suppression steepens the middle range, per-context damage saturation caps the small-$D$ end, and the replay-free residual bounds the large-$D$ end, where the memory-specific component of the gap vanishes. The function $a(\cdot)$ is an empirical factor; its microscopic origin (backbone readout calibration versus row-level traffic competition) remains open.

%Hash collisions may reduce the independence of local memory signals, but they can also mix unrelated contexts; they are one candidate micro-origin of the load factor above, not yet separated from the backbone-readout channel.

%\pending{Complete token-mass-weighted aggregation when the frequency distribution changes, the collision model, and matched data-volume experiments. The epoch-length figure is awaiting replacement. Closing the dose axis additionally requires the distinct-context count $K(D)$ at each dose.}

%\subsection{Further tests of multi-epoch dynamics}

%\pending{Derive the optimizer-dependent margin evolution and align it with measured per-pass damage. Distinguish approximately linear accumulation over a finite window, decreasing increments, and a possible long-run plateau. Specify observations that would falsify the model. The local margin, transfer coefficient, and aggregate gap slope currently serve as constraints, not a closed derivation.}

\subsection{Frequency-scaling fits}
\label{app:gt-fits}

The Good--Turing novel-continuation rate is recomputed from the raw training shard (seed-42 \code{shard_00001.bin}, 49.7M positions, chunked exactly as the data loader chunks it) as $\hat f(F)=N_1(F)/(F\,n(F))$, where $n(F)$ is the number of contexts with exact count $F$ and $N_1(F)$ the number of context--continuation types occurring once among them. Figure~\ref{fig:b_t_MF} shows every positive exact-$F$ layer. Beyond $F\approx10^3$ each layer holds only a handful of contexts (0.1\% of bigram contexts and fewer than 0.01\% of trigram contexts lie there) and the estimates fan out by two orders of magnitude from sampling noise. A fit that weights every layer equally over the full range is dominated by this tail and returns $\alpha=0.45$ (bigram) and $0.43$ (trigram); the version of this figure in an earlier draft, which fitted the per-count total $F\hat f(F)$ the same way, gave $0.34$ and $0.38$. We therefore fit on $F\le10^3$, the window in which the gap is measured and which holds 76\% (bigram) and 97\% (trigram) of training tokens, and weight each context equally, which is the estimator of the per-context mean that Equation~\eqref{eq:gap_j} averages. Under this weighting the tail no longer matters (the fit is unchanged if the window is widened to all $F$), and the exponents are $0.287$ and $0.314$. Weighting by tokens ($F\,n(F)$) instead gives $0.266$ and $0.246$; the spread between the two weightings, about $0.02$ for bigram and $0.07$ for trigram, is the corpus-side uncertainty of the exponent.

\begin{figure}[!htbp]
\centering
\includegraphics[width=\linewidth]{fig_app_gt_missing_mass.pdf}
\caption{Good--Turing novel-continuation rate $\hat f(F)$ by exact context count for bigram (left) and trigram (right) contexts. Coloured points are layers with at least 8 contexts, grey points sparser layers; the shaded region $F>10^3$ is excluded from the per-context-weighted fit (solid, extrapolated dashed). The dotted line is the equal-weight full-range fit.}
\label{fig:b_t_MF}
\end{figure}

Figure~\ref{fig:gt-kernel} compares the gap with $D\,\hat f(F)$ in the seven geometric frequency bins of the registered gap fit (run \code{s1v5_128_frequency_main}, seed 42, step 1000, fixed train probe and fixed validation loss, bins restricted to at least 32 shared contexts), pooling $\hat f$ over the same bins. With $D$ the geometric mean of $\operatorname{gap}/\hat f$ over the bins, the ratio is constant to within 9\% (bigram) and 18\% (trigram) across 2.5 decades of $F$, and $\log\operatorname{gap}$ correlates with $\log\hat f$ at $0.99$ and $0.98$. The trigram ratio drifts downward with $F$: the gap decays somewhat faster than the novel-continuation rate, so the depth per unit of novel mass, $D$, is not exactly frequency-independent for trigram contexts---rarer contexts are hurt slightly more per novel token, consistent with sharper memorization when fewer distinct continuations were seen. This, together with the 27\% of excess loss that falls on seen continuations of frequent contexts (Table~\ref{tab:novel-share}), is the part of the gap that coverage alone does not account for.

\begin{figure}[!htbp]
\centering
\includegraphics[width=\linewidth]{fig_app_gt_kernel.pdf}
\caption{Step-1000 gap against $D\,\hat f(F)$ in the seven matched frequency bins (top) and the per-bin ratio $\operatorname{gap}/\hat f$ with its mean and $\pm$CV band (bottom); dashed lines are the direct power-law fits of Equation~\eqref{eq:gap-fits}.}
\label{fig:gt-kernel}
\end{figure}
